% Standalone research note for the hybrid-local-solver project.
%
% This is a curated reconstruction of the project conversation on a
% center-star lower bound for AESP-PPR with the batched LocGD inner solver.
% It is not a verbatim transcript.  Incorrect intermediate claims are
% retained only as explicitly refuted or conditional proof routes.

\documentclass[11pt]{article}

\usepackage[margin=1in]{geometry}
\usepackage{amsmath,amssymb,amsthm,mathtools,bm}
\usepackage{booktabs,tabularx}
\usepackage{enumitem}
\usepackage{microtype}
\usepackage[numbers,sort&compress]{natbib}
\usepackage{xcolor}
\usepackage[colorlinks=true,citecolor=blue,linkcolor=black,urlcolor=blue]{hyperref}
\usepackage[capitalize,nameinlink,noabbrev]{cleveref}

\newtheorem{theorem}{Theorem}[section]
\newtheorem{lemma}[theorem]{Lemma}
\newtheorem{proposition}[theorem]{Proposition}
\newtheorem{corollary}[theorem]{Corollary}
\newtheorem{conjecture}[theorem]{Conjecture}
\theoremstyle{definition}
\newtheorem{definition}[theorem]{Definition}
\newtheorem{remark}[theorem]{Remark}

\newcommand{\R}{\mathbb{R}}
\newcommand{\one}{\mathbf{1}}
\newcommand{\e}{\mathrm{e}}
\newcommand{\norm}[1]{\left\lVert #1\right\rVert}
\newcommand{\abs}[1]{\left\lvert #1\right\rvert}
\newcommand{\vol}{\operatorname{vol}}
\newcommand{\supp}{\operatorname{supp}}
\newcommand{\Work}{\mathcal{W}}
\newcommand{\AESP}{\textnormal{\textsc{AESP}}}
\newcommand{\LOCGD}{\textnormal{\textsc{LocGD}}}
\newcommand{\LOCAPPR}{\textnormal{\textsc{LocAPPR}}}
\newcommand{\APPR}{\textnormal{\textsc{APPR}}}
\newcommand{\softO}{\widetilde{\mathcal O}}
\newcommand{\statussource}{\textbf{Source result}}
\newcommand{\statusproved}{\textbf{Proved here}}
\newcommand{\statusconditional}{\textbf{Conditional}}
\newcommand{\statusrefuted}{\textbf{Refuted route}}
\newcommand{\statusopen}{\textbf{Open}}

\title{A Star-Graph Work Lower Bound for AESP with Batched Local Gradient Descent}
\author{Baojian Zhou\\
\small Working research note for \texttt{hybrid-local-solver}}
\date{August 2026}

\begin{document}
\maketitle

\begin{abstract}
We prove a worst-case lower bound for the literal personalized PageRank
specialization of the accelerated evolving-set process (AESP) when every
proximal subproblem is solved by the batched local gradient method LocGD.  On
a center-seeded star with $B$ leaves, every nonempty inner batch scans degree
volume at least $B$.  The nontrivial part is to show that Catalyst momentum
cannot turn only $o(1/\sqrt{\alpha})$ such batches into a constant amount of
PageRank mass.  A direct residual-cone argument is only conditional, because
outer extrapolation need not preserve the cone.  We close this gap using the
first nonempty call: its activation condition forces all subsequent AESP
subproblem tolerances to be on the $\alpha^2/B$ scale.  The certified
subproblem objective gap then bounds every slow-mode proximal error by
$\sqrt{2}\alpha/(1-\alpha)$.  A positive Green-function calculation controls
all momentum amplification and yields
\[
  \Work_{\AESP\text{-}\LOCGD}
  =\Omega\!\left(\frac{1}{\sqrt{\alpha}\,\epsilon}\right).
\]
More generally, the worst-case work over graphs with at most $m$ edges is
$\Omega(\min\{m,1/\epsilon\}/\sqrt{\alpha})$.  The theorem is unconditional
within its stated algorithm and work model; it is not a lower bound for an
arbitrary AESP inner map, sequential LocAPPR, or every hybrid local solver.
\end{abstract}

\tableofcontents

\section{Scope, status, and source map}
\label{sec:scope}

This note reconstructs the complete proof development leading to the
center-star lower bound.  It deliberately separates the literal source
algorithm from new arguments and from proof paths that did not close.

\begin{center}
\begin{tabularx}{\textwidth}{@{}p{0.19\textwidth}X@{}}
\toprule
Status & Statement \\
\midrule
\statussource & The AESP-PPR parameters, LocGD activation threshold,
subproblem-gap certificate, stopping rule, and active-volume work model are
those of Algorithms 2--3 and Lemma 3.2 of \citet{huang2025aesp}. \\
\statusproved & Batched LocGD preserves leaf symmetry; every nonempty star
batch costs at least $B$; and the literal AESP-LocGD algorithm has work
$\Omega(B/\sqrt{\alpha})$ whenever $B\epsilon\leq 1/4$. \\
\statusproved & Choosing $B=\lfloor 1/(4\epsilon)\rfloor$ gives
$\Omega(1/(\sqrt{\alpha}\epsilon))$; imposing an edge budget gives
$\Omega(\min\{m,1/\epsilon\}/\sqrt{\alpha})$. \\
\statusconditional & A shorter slow-innovation proof works if a
center--leaf residual cone is invariant at every outer point. \\
\statusrefuted & Leaf symmetry alone does not imply that cone and does not by
itself imply an $O(\alpha)$ innovation bound. \\
\statusopen & A polynomially larger lower bound for AESP-LocGD, or the same
lower bound for arbitrary AESP inner operators and hybrid local methods,
requires a different argument. \\
\bottomrule
\end{tabularx}
\end{center}

The source locations used here are as follows.  In \citet{huang2025aesp}, PDF
pages 3--5 define the PPR quadratic, inexact proximal set, nested work, and
inner threshold; PDF pages 5--7 state the LocGD update and AESP-PPR parameters;
PDF pages 16--21 prove the inner gap and convergence claims; and PDF page 29
gives the batched queue implementation of Algorithm 3.  These locations are
recorded in the repository literature note as well.

\paragraph{Note-scoped convention.}
The repository-wide residual convention remains open.  Throughout this note,
$\epsilon$ means only the AESP-PPR output accuracy
\begin{equation}
  \norm{D^{-1}(\widehat\pi-\pi)}_\infty\leq\epsilon,
  \label{eq:accuracy}
\end{equation}
and $\Work$ means only the cumulative active degree volume in
\cref{def:work}.  No conversion to an RPPR proximal-gradient residual, APPR
threshold, or implementation-wide stopping rule is asserted.

\section{The literal AESP-LocGD model}
\label{sec:model}

Let $G=(V,E)$ be a connected undirected unweighted graph with adjacency
matrix $A$, degree matrix $D$, and $m=|E|$.  For a source distribution
$s\geq 0$ with $\one^\top s=1$, the symmetrized PPR quadratic is
\begin{equation}
  f(x)
  :=\frac12 x^\top Qx-\alpha x^\top D^{-1/2}s,
  \qquad
  Q:=\frac{1+\alpha}{2}I
     -\frac{1-\alpha}{2}D^{-1/2}AD^{-1/2}.
  \label{eq:ppr-objective}
\end{equation}
Its unique minimizer $x^\star$ satisfies $\pi=D^{1/2}x^\star$.
The eigenvalues of $Q$ lie in $[\alpha,1]$.

We use the PPR specialization of AESP for
\begin{equation}
  0<\alpha<\frac12,
  \qquad
  \eta:=1-2\alpha,
  \qquad
  d:=\alpha+\eta=1-\alpha.
  \label{eq:parameters}
\end{equation}
At outer iteration $t$, given the extrapolated point $y^{(t-1)}$, LocGD
approximately minimizes
\begin{equation}
  h_t(z):=f(z)+\frac{\eta}{2}\norm{z-y^{(t-1)}}_2^2
  \label{eq:subproblem}
\end{equation}
from $z_t^{(0)}=y^{(t-1)}$.  The required gap and its scheduled tolerance are
\begin{align}
  H_t(\phi_t)
  &:=\{z:h_t(z)-h_t^\star\leq\phi_t\},
  \label{eq:Ht}\\
  \phi_t
  &:=\frac{1+\alpha}{18}
      \left(1-\frac9{10}\sqrt{\frac{\alpha}{1-\alpha}}\right)^t.
  \label{eq:phi}
\end{align}
The inner normalized-gradient threshold in Algorithm 3 is
\begin{equation}
  \varepsilon_t^{\mathrm{in}}
  :=\max\left\{
       \sqrt{\frac{d\phi_t}{m}},
       \frac{2d\phi_t}
       {\norm{D^{1/2}\nabla h_t(z_t^{(0)})}_1}
     \right\}.
  \label{eq:inner-threshold}
\end{equation}
If the denominator vanishes, Algorithm 3 returns the exact initial minimizer.
Otherwise, a vertex is active when
\begin{equation}
  \abs{[\nabla h_t(z)]_v}
  \geq \varepsilon_t^{\mathrm{in}}\sqrt{d_v}.
  \label{eq:activation}
\end{equation}
All currently queued vertices form one batch and are updated with step size
\begin{equation}
  \tau_{\mathrm{gd}}:=\frac{2}{1+\alpha+2\eta}
  =\frac{2}{3(1-\alpha)}.
  \label{eq:gd-step}
\end{equation}
The source analysis guarantees that the returned point $x^{(t)}$ belongs to
$H_t(\phi_t)$, including when the queue is initially empty.  The outer update
is
\begin{equation}
  y^{(t)}=x^{(t)}+\beta(x^{(t)}-x^{(t-1)}),
  \qquad
  \beta:=\frac{\sqrt{1-\alpha}-\sqrt\alpha}
               {\sqrt{1-\alpha}+\sqrt\alpha}.
  \label{eq:momentum}
\end{equation}

\begin{definition}[Work measure]
\label{def:work}
If $S_t^{(k)}$ is the active set in inner batch $k$ of outer call $t$, define
\begin{equation}
  \Work:=\sum_t\sum_{k=0}^{K_t-1}\vol(S_t^{(k)}),
  \qquad
  \vol(S):=\sum_{v\in S}d_v.
  \label{eq:work}
\end{equation}
This is exactly $\sum_t K_t\overline{\vol}(S_t)$ in the AESP paper.  It does
not charge an empty outer call.  A separate augmented measure that charges
outer-loop overhead is considered in \cref{sec:refinements}.
\end{definition}

\section{Center-star reduction}
\label{sec:star}

Let $G_B=K_{1,B}$, let $c$ denote its center, let $L$ denote its $B$ leaves,
and take the source $s=e_c$.  The center has degree $B$ and every leaf has
degree one.  All iterates remain leaf-symmetric, so write $x_\ell$ for the
common leaf coordinate and define
\begin{equation}
  C(x):=\sqrt{B}\,x_c,
  \qquad
  L(x):=B x_\ell,
  \qquad
  M(x):=C(x)+L(x),
  \qquad
  F(x):=C(x)-L(x).
  \label{eq:mass-coordinates}
\end{equation}
Here $C$ is the center PageRank mass, $L$ is aggregate leaf mass, $M$ is the
slow stationary mode, and $F$ is the fast alternating mode.

\begin{lemma}[Exact solution and norm conversion]
\label{lem:norm-conversion}
On the center-seeded star,
\begin{equation}
  C^\star=\frac{1+\alpha}{2},
  \qquad
  L^\star=\frac{1-\alpha}{2},
  \qquad
  M^\star=1,
  \qquad
  F^\star=\alpha.
  \label{eq:star-solution}
\end{equation}
For every leaf-symmetric $x$ and $\widehat\pi=D^{1/2}x$,
\begin{equation}
  \norm{D^{-1}(\widehat\pi-\pi)}_\infty
  =\frac{
      \abs{M(x)-1}+\abs{F(x)-\alpha}
    }{2B}.
  \label{eq:star-error}
\end{equation}
Consequently, if $B\epsilon\leq 1/4$, every output satisfying
\cref{eq:accuracy} obeys $M(x)\geq 1/2$.
\end{lemma}

\begin{proof}
The PPR fixed point reduces to the two-state lazy star walk, giving
\cref{eq:star-solution}.  The center and per-leaf errors are
\begin{align*}
  C-C^\star
  &=\frac{(M-1)+(F-\alpha)}2,\\
  \frac{L-L^\star}{B}
  &=\frac{(M-1)-(F-\alpha)}{2B}.
\end{align*}
After degree normalization both coordinates have denominator $B$.  The
identity $\max\{|a+b|,|a-b|\}=|a|+|b|$ proves
\cref{eq:star-error}.  Accuracy therefore implies
$|M-1|\leq2B\epsilon\leq1/2$.
\end{proof}

\begin{lemma}[Symmetry and scan cost]
\label{lem:scan-cost}
For AESP with batched LocGD on $G_B$, every inner active set is one of
\begin{equation}
  \varnothing,\qquad \{c\},\qquad L,\qquad \{c\}\cup L.
  \label{eq:active-sets}
\end{equation}
Every nonempty inner batch therefore has volume at least $B$.  In particular,
if $K$ outer calls contain at least one nonempty batch, then
\begin{equation}
  \Work\geq BK.
  \label{eq:work-vs-calls}
\end{equation}
\end{lemma}

\begin{proof}
The zero initialization is invariant under every permutation of the leaves.
The AESP extrapolation is linear and permutation-equivariant.  At an inner
point with identical leaves, their gradients and activation tests are
identical.  A center update changes all leaf gradients by the same amount,
and a simultaneous leaf update treats all leaves identically.  Induction over
the outer calls and inner batches proves \cref{eq:active-sets}.  Finally,
\[
  \vol(\{c\})=B,
  \qquad
  \vol(L)=B,
  \qquad
  \vol(V)=2B.
\]
\end{proof}

Thus the lower bound reduces to showing that at least
$\Omega(1/\sqrt\alpha)$ outer calls must be effective.

\section{The residual-cone route and why it is insufficient}
\label{sec:cone}

This section records the first proof route.  Its inner-process algebra is
useful, but its proposed outer invariant is not implied by symmetry.

For a fixed subproblem, let $r_c$ and $r_L$ be the negative-gradient residual
in the mass coordinates.  The shifted Hessian on the $(C,L)$ subspace is
\begin{equation}
  H=Q+\eta I
   =\frac{1-\alpha}{2}
     \begin{pmatrix}3&-1\\-1&3\end{pmatrix}.
  \label{eq:shifted-hessian}
\end{equation}
Using \cref{eq:gd-step}, one LocGD batch transforms the residual exactly as
\begin{equation}
\begin{array}{c|cc}
\text{active block}&r_c^+&r_L^+\\
\hline
\{c\}&0&r_L+r_c/3\\
L&r_c+r_L/3&0\\
V&r_L/3&r_c/3.
\end{array}
\label{eq:residual-map}
\end{equation}
At the initial point $z=y$ of an inner call, put
\begin{equation}
  \sigma:=r_c+r_L=\alpha(1-M(y)),
  \qquad
  \Delta:=r_c-r_L=\alpha-F(y).
  \label{eq:sigma-delta}
\end{equation}

\begin{lemma}[Conditional slow innovation]
\label{lem:conditional-innovation}
Suppose an inner call begins in the cone
\begin{equation}
  \sigma\geq0,
  \qquad
  |\Delta|\leq2\sigma.
  \label{eq:cone}
\end{equation}
If $x$ is the LocGD output and $u:=M(x)-M(y)$, then
\begin{equation}
  0\leq u
  \leq\frac{\alpha}{1-\alpha}(1-M(y)).
  \label{eq:cone-innovation}
\end{equation}
\end{lemma}

\begin{proof}
If $r_c,r_L\geq0$, every map in \cref{eq:residual-map} preserves
nonnegativity and does not increase their sum.  Otherwise, assume without
loss of generality that $r_c=-a<0$ and $r_L=b>0$.  The cone gives
$a+b\leq2(b-a)$, hence $b\geq3a$.  Because the activation threshold is common
in mass coordinates, the negative coordinate cannot be active without the
positive one.  If only the positive leaf block is active, the new residual is
$(-a+b/3,0)\geq0$.  If both blocks are active, it is
$(b/3,-a/3)$ and the same dominance is preserved.  Thus the residual sum
remains nonnegative and nonincreasing throughout the inner process.

The slow eigenvalue of \cref{eq:shifted-hessian} is $1-\alpha$, so
\[
  \sigma_{\mathrm{out}}
  =\sigma_{\mathrm{in}}-(1-\alpha)u.
\]
Combining $0\leq\sigma_{\mathrm{out}}\leq\sigma_{\mathrm{in}}$ with
\cref{eq:sigma-delta} proves the result.
\end{proof}

The missing step would have been to prove \cref{eq:cone} at every AESP outer
point.  That claim is stronger than symmetry and is not automatic under
extrapolation.  For example, a symmetric point with
\[
  M=\frac12,
  \qquad
  F=\frac14
\]
has $\sigma=\alpha/2$ and $\Delta=\alpha-1/4$.  For small $\alpha$ the cone
fails by a factor of order $1/\alpha$.  A common threshold can then select the
positive coordinate while leaving a negative coordinate untouched, creating
an order-one slow innovation.  This example need not itself be an AESP
iterate; it is enough to refute the inference ``leaf symmetry implies the
cone.''  The unconditional proof must use the actual tolerance schedule.

\section{Unconditional lower bound for literal AESP-LocGD}
\label{sec:main-proof}

Set
\begin{equation}
  q:=\frac{\alpha}{1-\alpha},
  \qquad
  \delta:=\sqrt q,
  \qquad
  \beta=\frac{1-\delta}{1+\delta},
  \qquad
  r:=1-\delta.
  \label{eq:q-delta}
\end{equation}
Because $\alpha<1/2$, we have $0<\delta<1$ and $0<\beta<1$.

\subsection{The first nonempty call fixes the error scale}

At the start of an inner call, define the symmetric residual masses
\begin{equation}
  r_c:=-\sqrt B\,[\nabla h_t(y^{(t-1)})]_c,
  \qquad
  r_L:=-B[\nabla h_t(y^{(t-1)})]_\ell,
  \qquad
  R_t:=|r_c|+|r_L|.
  \label{eq:mass-residuals}
\end{equation}
Symmetry gives
$R_t=\norm{D^{1/2}\nabla h_t(y^{(t-1)})}_1$.  Since $m=B$, multiplying
\cref{eq:inner-threshold} by $B$ shows that the common activation threshold in
the mass coordinates is
\begin{equation}
  \Theta_t
  :=B\varepsilon_t^{\mathrm{in}}
  =\max\left\{A_t,\frac{2A_t^2}{R_t}\right\},
  \qquad
  A_t:=\sqrt{B(1-\alpha)\phi_t}.
  \label{eq:mass-threshold}
\end{equation}

\begin{lemma}[First-activation certificate]
\label{lem:first-activation}
Let $t_0$ be the first outer call with a nonempty LocGD batch.  Then
\begin{equation}
  A_{t_0}\leq\frac{\alpha}{\sqrt2},
  \qquad
  A_t\leq\frac{\alpha}{\sqrt2}
  \quad\text{for all }t\geq t_0.
  \label{eq:A-bound}
\end{equation}
\end{lemma}

\begin{proof}
Every earlier empty call returns its initial point.  Starting from
$x^{(0)}=y^{(0)}=0$, all points before $t_0$ therefore remain zero.  At
$t_0$,
\[
  r_c=\alpha,
  \qquad
  r_L=0,
  \qquad
  R_{t_0}=\alpha.
\]
The center is active only if
\[
  \alpha\geq\Theta_{t_0}
  =\max\left\{A_{t_0},\frac{2A_{t_0}^2}{\alpha}\right\}.
\]
The second term gives $A_{t_0}\leq\alpha/\sqrt2$.  The schedule
\cref{eq:phi} decreases with $t$, so the same bound holds later.
\end{proof}

\subsection{Certified slow-mode accuracy of every later inner output}

Let
\begin{equation}
  M_t:=M(x^{(t)}),
  \qquad
  Y_{t-1}:=M(y^{(t-1)}).
  \label{eq:M-Y}
\end{equation}
The normalized stationary eigenvector is
\[
  w:=\frac{D^{1/2}\one}{\sqrt{2B}},
  \qquad
  Qw=\alpha w.
\]
Multiplying the exact proximal system
\[
  (Q+\eta I)\widetilde x_t
  =\alpha D^{-1/2}e_c+\eta y^{(t-1)}
\]
by $\one^\top D^{1/2}$ gives its exact slow mass
\begin{equation}
  \widetilde M_t
  =\frac{\alpha+\eta Y_{t-1}}{\alpha+\eta}
  =q+(1-q)Y_{t-1}.
  \label{eq:exact-prox-mass}
\end{equation}

\begin{lemma}[Uniform slow-mode proximal error]
\label{lem:prox-error}
Renumber $t_0$ as outer call $1$.  For every subsequent returned point,
including an empty call,
\begin{equation}
  M_t=q+(1-q)Y_{t-1}+\xi_t,
  \qquad
  |\xi_t|\leq\sqrt2\,q.
  \label{eq:xi-bound}
\end{equation}
\end{lemma}

\begin{proof}
The source guarantee gives $x^{(t)}\in H_t(\phi_t)$.  Since the slow
eigenvalue of $Q+\eta I$ is $1-\alpha$, strong convexity projected onto $w$
gives
\begin{equation}
  h_t(x^{(t)})-h_t^\star
  \geq\frac{1-\alpha}{4B}
       (M_t-\widetilde M_t)^2.
  \label{eq:gap-projection}
\end{equation}
For completeness, the full two-mode decomposition is
\[
  h_t(x)-h_t^\star
  =\frac{1-\alpha}{4B}(M(x)-\widetilde M_t)^2
   +\frac{1-\alpha}{2B}(F(x)-\widetilde F_t)^2.
\]
Combining \cref{eq:gap-projection} with the gap certificate and
\cref{eq:A-bound},
\[
  |M_t-\widetilde M_t|
  \leq2\sqrt{\frac{B\phi_t}{1-\alpha}}
  =\frac{2A_t}{1-\alpha}
  \leq\sqrt2\frac{\alpha}{1-\alpha}
  =\sqrt2\,q.
\]
This proves \cref{eq:xi-bound}.  If the initial gradient is zero, the output
is exact; if the queue is merely empty, Lemma 3.2 of the source still certifies
the displayed objective gap.
\end{proof}

\subsection{Positive Green function and momentum control}

Using $Y_{t-1}=(1+\beta)M_{t-1}-\beta M_{t-2}$, the slow recurrence is
\begin{equation}
  M_t
  =q+(1-q)\bigl((1+\beta)M_{t-1}-\beta M_{t-2}\bigr)+\xi_t.
  \label{eq:slow-recurrence}
\end{equation}
All pre-$t_0$ iterates vanish, so after renumbering
$M_{-1}=M_0=0$.  The identities
\begin{equation}
  (1-q)(1+\beta)=2r,
  \qquad
  (1-q)\beta=r^2
  \label{eq:critical-identities}
\end{equation}
turn \cref{eq:slow-recurrence} into
\begin{equation}
  M_t=2rM_{t-1}-r^2M_{t-2}+q+\xi_t.
  \label{eq:critical-recurrence}
\end{equation}

Let $\overline M_t$ denote the exact-proximal trajectory, obtained by setting
$\xi_t=0$.  Solving the repeated-root recurrence gives
\begin{equation}
  \overline M_t
  =1-(1+t\delta)(1-\delta)^t.
  \label{eq:exact-trajectory}
\end{equation}
Its extrapolated mass is
\begin{equation}
  \overline Y_t
  =(1+\beta)\overline M_t-\beta\overline M_{t-1},
  \qquad
  1-\overline Y_t
  =\frac{1+(t+1)\delta}{1+\delta}(1-\delta)^t.
  \label{eq:exact-Y}
\end{equation}
Hence $0\leq\overline Y_t\leq1$.

\begin{lemma}[Positive impulse response]
\label{lem:green}
The impulse response from $\xi_j$ to $M_{j+k}$ in
\cref{eq:critical-recurrence} is
\begin{equation}
  h_k=(k+1)r^k,
  \qquad
  \sum_{k=0}^\infty h_k=\frac1q.
  \label{eq:h-green}
\end{equation}
The impulse response from $\xi_j$ to the corresponding extrapolated mass is
nonnegative and equals
\begin{equation}
  g_0=1+\beta,
  \qquad
  g_k=(1+\beta)h_k-\beta h_{k-1}
     =\beta(k+2)r^{k-1}\quad(k\geq1),
  \qquad
  \sum_{k=0}^\infty g_k=\frac1q.
  \label{eq:g-green}
\end{equation}
Consequently,
\begin{equation}
  |Y_{t-1}-\overline Y_{t-1}|\leq\sqrt2,
  \qquad
  Y_{t-1}\geq-\sqrt2.
  \label{eq:Y-lower}
\end{equation}
\end{lemma}

\begin{proof}
The transfer polynomial of the homogeneous recurrence is
$(z-r)^2$, giving $h_k=(k+1)r^k$.  Summing the geometric derivative gives
$\sum_k h_k=(1-r)^{-2}=\delta^{-2}=q^{-1}$.  Applying the extrapolation
filter $(1+\beta)-\beta z^{-1}$ gives \cref{eq:g-green}; direct substitution
using $\beta=r/(1+\delta)$ shows every coefficient is nonnegative.  Its sum is
the sum of $h_k$ because the filter has unit gain at one.  Therefore
\[
  |Y_{t-1}-\overline Y_{t-1}|
  \leq \sup_j|\xi_j|\sum_{k\geq0}g_k
  \leq(\sqrt2 q)q^{-1}=\sqrt2.
\]
Now use $\overline Y_{t-1}\geq0$.
\end{proof}

\subsection{Innovation counting}

Define the innovation created by outer call $t$ as
\begin{equation}
  u_t:=M_t-Y_{t-1}.
  \label{eq:innovation}
\end{equation}
The identities \cref{eq:exact-prox-mass,eq:xi-bound,eq:Y-lower} imply
\begin{equation}
  u_t=q(1-Y_{t-1})+\xi_t
  \leq(1+2\sqrt2)q.
  \label{eq:uniform-innovation}
\end{equation}
If call $t$ is empty, its output equals its initialization, so $u_t=0$.

Independently of the proximal recurrence, the momentum identity gives
\begin{equation}
  M_t=(1+\beta)M_{t-1}-\beta M_{t-2}+u_t.
  \label{eq:innovation-recurrence}
\end{equation}
Starting from zero and unrolling,
\begin{equation}
  M_T
  =\sum_{t=1}^T
     u_t\frac{1-\beta^{T-t+1}}{1-\beta}.
  \label{eq:innovation-unroll}
\end{equation}
All weights are positive and at most $(1-\beta)^{-1}$.  Negative innovations
only lower the right-hand side, so if $K$ calls are effective,
\begin{equation}
  M_T\leq
  \frac{(1+2\sqrt2)qK}{1-\beta}.
  \label{eq:M-upper}
\end{equation}

\begin{theorem}[Center-star lower bound]
\label{thm:star-lower-bound}
Let $0<\alpha<1/2$, let $G_B=K_{1,B}$ have its center as source, and run the
literal AESP-PPR algorithm with the batched LocGD inner solver from zero.  If
$B\epsilon\leq1/4$, then every output satisfying \cref{eq:accuracy} requires
\begin{equation}
  \boxed{
  \Work
  \geq
  \frac{B}
  {(1+2\sqrt2)\delta(1+\delta)},
  \qquad
  \delta=\sqrt{\frac{\alpha}{1-\alpha}}.}
  \label{eq:main-lower-bound}
\end{equation}
In particular,
\begin{equation}
  \Work=\Omega\!\left(B\sqrt{\frac{1-\alpha}{\alpha}}\right).
  \label{eq:B-alpha-asymptotic}
\end{equation}
\end{theorem}

\begin{proof}
By \cref{lem:norm-conversion}, accuracy forces $M_T\geq1/2$.  Combining this
with \cref{eq:M-upper} gives
\[
  K\geq\frac{1-\beta}{2(1+2\sqrt2)q}.
\]
Since
\[
  1-\beta=\frac{2\delta}{1+\delta},
  \qquad q=\delta^2,
\]
we obtain
\[
  K\geq\frac1{(1+2\sqrt2)\delta(1+\delta)}.
\]
Apply \cref{lem:scan-cost}.
\end{proof}

\begin{corollary}[Accuracy-parametrized lower bound]
\label{cor:epsilon-bound}
For $0<\epsilon\leq1/8$, choose
$B=\lfloor1/(4\epsilon)\rfloor$.  Then
\begin{equation}
  \boxed{
  \Work_{\AESP\text{-}\LOCGD}
  \geq
  \frac{1}{16\sqrt2(1+2\sqrt2)}
  \frac1{\sqrt\alpha\,\epsilon}.}
  \label{eq:epsilon-lower-bound}
\end{equation}
\end{corollary}

\begin{proof}
For $\epsilon\leq1/8$,
$B\geq1/(8\epsilon)$.  Also $\delta<1$, so
$(1+\delta)^{-1}>1/2$, and $1-\alpha>1/2$, so
$\delta^{-1}\geq1/\sqrt{2\alpha}$.  Substitute these estimates into
\cref{eq:main-lower-bound}.
\end{proof}

\section{Exact-proximal outer obstruction}
\label{sec:exact-prox}

The preceding theorem concerns the literal thresholded inner solver.  A
simpler companion calculation isolates the outer slow-mode obstruction.  For
exact proximal solves, \cref{eq:exact-trajectory} gives the normalized slow
error
\begin{equation}
  E_t:=1-\overline M_t
  =(1+t\delta)(1-\delta)^t.
  \label{eq:exact-error}
\end{equation}
If $t\delta\leq1/4$, then Bernoulli's inequality yields
\[
  E_t\geq(1-\delta)^t\geq1-t\delta\geq\frac34,
\]
so $\overline M_t\leq1/4$.  Hence even exact-proximal AESP needs
$\Omega(1/\sqrt\alpha)$ outer calls to create constant slow mass.  This
calculation is not itself a work lower bound for LocGD: the literal theorem
additionally needs \cref{lem:first-activation,lem:prox-error,lem:green} to
connect each effective call to degree work.

\section{Refinements and prospects for stronger lower bounds}
\label{sec:refinements}

\subsection{A graph-budget statement}

\begin{corollary}[Worst case under an edge budget]
\label{cor:edge-budget}
Let $m\geq1$ and $0<\epsilon\leq1/8$.  Among connected graphs with at most
$m$ edges, AESP-PPR with batched LocGD has worst-case work
\begin{equation}
  \boxed{
  \Omega\!\left(
    \frac{\min\{m,1/\epsilon\}}{\sqrt\alpha}
  \right).}
  \label{eq:edge-budget}
\end{equation}
\end{corollary}

\begin{proof}
Choose the star with
$B=\min\{m,\lfloor1/(4\epsilon)\rfloor\}$.  It has at most $m$ edges and
$B\epsilon\leq1/4$.  Moreover,
$B=\Omega(\min\{m,1/\epsilon\})$.  Apply
\cref{thm:star-lower-bound}.
\end{proof}

When $m=\Theta(1/\epsilon)$, the construction forces
$\Omega(m/\sqrt\alpha)$ work: every effective round scans essentially the
whole graph.

\subsection{Outer-loop overhead ignored by active volume}

The work measure \cref{eq:work} charges an empty inner call zero, although the
implementation still initializes its tolerance and evaluates its queue.  The
first-activation certificate gives a separate delay.  From
\cref{eq:phi,eq:mass-threshold},
\[
  A_t^2
  =\frac{B(1-\alpha^2)}{18}
    \left(1-\frac9{10}\delta\right)^t.
\]
The necessary condition $A_{t_0}^2\leq\alpha^2/2$ therefore implies
\begin{equation}
  t_0\geq
  \frac{
    \log\!\left(B(1-\alpha^2)/(9\alpha^2)\right)
  }{
    -\log(1-0.9\delta)
  }
  \label{eq:first-call-delay}
\end{equation}
whenever the numerator is positive.  If one defines an augmented cost
$\Work_+:=\Work+T$ that charges one unit per outer call, then for small
$\alpha$,
\begin{equation}
  \Work_+
  =\Omega\!\left(
     \frac{B}{\sqrt\alpha}
     +\frac1{\sqrt\alpha}
      \log_+\frac{B}{\alpha^2}
   \right).
  \label{eq:augmented-work}
\end{equation}
The logarithmic term is genuine implementation overhead but is not part of
the dominated active-volume measure used in the AESP theorem.

\subsection{Why the star does not currently give a polynomially larger bound}

The proved result does not establish either
\[
  \Omega\!\left(\frac1{\alpha\epsilon}\right)
  \qquad\text{or}\qquad
  \Omega\!\left(\frac1{\sqrt\alpha\,\epsilon^2}\right)
\]
for AESP-LocGD.  The star exposes two obstructions and no more:
\begin{enumerate}[leftmargin=*]
\item its stationary slow mode needs $\Theta(1/\sqrt\alpha)$ accelerated
      outer progress; and
\item every effective symmetric scan costs $\Theta(B)$.
\end{enumerate}
Accuracy forces constant mass only while $B=O(1/\epsilon)$, so this direct
product saturates at $1/(\sqrt\alpha\epsilon)$.  Even a hypothetical
precision-sensitive refinement
$B\log(1/(B\epsilon))/\sqrt\alpha$ would optimize to only
$\Theta(1/(\sqrt\alpha\epsilon))$ over $B\epsilon\leq1/4$.

A stronger active-volume lower bound therefore needs a graph that forces
transient exploration beyond the $O(1/\epsilon)$ volume needed by the final
significant output.  Plausible candidates are a multiscale spider, a
star-of-stars, or a clustered lollipop in which momentum repeatedly crosses
successive activation thresholds.  A successful construction must prove all
of the following simultaneously:
\begin{enumerate}[leftmargin=*]
\item the high-volume blocks are actually activated by the literal AESP
      threshold, not merely by an unconstrained accelerated recurrence;
\item cancellation or signed residuals cannot bypass the intended scans;
\item the required output accuracy remains nontrivial under degree
      normalization; and
\item the total forced work is not already limited by a matching explicit
      AESP trajectory.
\end{enumerate}
At present these are research directions, not lower-bound claims.

\section{Proof audit and exact scope}
\label{sec:audit}

The proof depends on the following chain, each link of which is explicit:
\begin{enumerate}[leftmargin=*]
\item AESP-PPR uses $\eta=1-2\alpha$, the schedule \cref{eq:phi}, and the
      LocGD threshold \cref{eq:inner-threshold}.
\item Batched queue processing, rather than sequential LocAPPR, preserves the
      all-leaf active block and forces cost at least $B$.
\item The first nonempty call starts from zero and certifies
      $A_{t_0}\leq\alpha/\sqrt2$.
\item The source inner guarantee $x^{(t)}\in H_t(\phi_t)$ converts that
      certificate into the uniform slow error $|\xi_t|\leq\sqrt2q$.
\item The critically damped outer recurrence has nonnegative impulse
      coefficients whose total gain is exactly $1/q$.
\item Therefore each effective call contributes at most $O(q)$ innovation
      before a total momentum amplification of $O(1/\sqrt q)$.
\item Accuracy forces constant slow mass, hence
      $\Omega(1/\sqrt q)=\Omega(1/\sqrt\alpha)$ effective calls.
\item Every such call costs at least $B$, and $B=\Theta(1/\epsilon)$ is legal.
\end{enumerate}

No step invokes the residual cone in \cref{sec:cone}.  Thus
\cref{thm:star-lower-bound} is unconditional for the stated algorithm.  The
following extensions are \emph{not} proved by this note:
\begin{itemize}[leftmargin=*]
\item replacing batched LocGD by sequential LocAPPR or an arbitrary map $M$;
\item changing the AESP tolerance schedule, momentum, or initialization;
\item transferring the result to $\ell_1$-regularized PageRank;
\item lower-bounding every accelerated local first-order or hybrid method; or
\item proving that $1/(\sqrt\alpha\epsilon)$ is optimal for the entire local
      oracle model.
\end{itemize}

\section{Conclusion}

The center-star argument rigorously matches the hoped-for hybrid target
$1/(\sqrt\alpha\epsilon)$ for one important baseline: literal AESP-PPR with
batched LocGD.  The key correction to the initial proof attempt is that
symmetry controls scan cost but not signed slow-mode progress.  The actual
AESP threshold supplies the missing control: the first effective call forces
all later proximal errors to the $O(\alpha)$ scale, and the positive Green
function bounds their total momentum amplification.  The resulting lower
bound is publication-ready at its stated scope, while a stronger or
algorithm-independent lower bound remains open.

\footnotesize
\begin{thebibliography}{9}

\bibitem[Huang et~al.(2025)Huang, Luo, Xiao, Yang, and Zhou]{huang2025aesp}
Binbin Huang, Luo Luo, Yanghua Xiao, Deqing Yang, and Baojian Zhou.
\newblock Accelerated evolving set processes for local PageRank computation.
\newblock In \emph{Advances in Neural Information Processing Systems}, 2025.

\bibitem[Zhou et~al.(2024)Zhou, Sun, Babanezhad Harikandeh, Guo, Yang, and
Xiao]{zhou2024locesp}
Baojian Zhou, Yifan Sun, Reza Babanezhad Harikandeh, Xingzhi Guo, Deqing Yang,
and Yanghua Xiao.
\newblock Iterative methods via locally evolving set process.
\newblock In \emph{Advances in Neural Information Processing Systems}, 2024.

\bibitem[Lin et~al.(2018)Lin, Mairal, and Harchaoui]{lin2018catalyst}
Hongzhou Lin, Julien Mairal, and Zaid Harchaoui.
\newblock Catalyst acceleration for first-order convex optimization: From
theory to practice.
\newblock \emph{Journal of Machine Learning Research}, 18(212):1--54, 2018.

\end{thebibliography}

\end{document}
